% Standalone section. Paste into the weekly report, or \input it.
% Week 4 note: the locked panel of geometric measures this thesis extracts from a centerline,
% written for a reader meeting the concept for the first time, because the two existing notes it
% draws on (thesis/week2/hemodynamics.tex and thesis/week3/tortuosity_and_representation.tex) were
% written for a supervisor who already knows the field.
%
% READING DISCIPLINE. No new reading was done for this note; every number and every citation below
% reuses a fact already recorded, with its own provenance, in the two notes above and in
% thesis/clinical_background/05_hemodynamics.tex:
%   - kashyap2022tortuosity: R^2 and p-values read from the full text (PMC8764056, Tables 1 and 4),
%     per week3's own disclosure.
%   - telloayala2026tortuosity, hart1999tortuosity, grisan2008tortuosity, bullitt2003tortuosity,
%     groves2009tortuosity: read at abstract level only, per week3's own disclosure. Re-read before
%     this note is merged into the thesis proper.
%   - asakura1990flow, chatzizisis2007ess, zhang2024curvature, li2011tortuosity, griffo2026wss,
%     taylor2024murray: as read and cited in thesis/clinical_background/05_hemodynamics.tex and
%     thesis/week2/hemodynamics.tex.
%   - shen2026geometry is cited for framing only, second-hand, exactly as week2 and week3 flag it.
%
% RECONCILIATION WITH THE DEEP-RESEARCH PASS. An earlier pass at this same question, run without
% access to this repository's own reading, proposed the Horton-Strahler bifurcation ratio and
% torsion as the two locked "advanced" methods. That pass had no primary-source WSS evidence for
% either. This note supersedes it: the locked panel below follows the evidence and the reasoning
% already worked out in week2 and week3, which is stronger because it is anchored in numbers read
% from the primary text rather than search summaries.
%
% Citation keys resolve against thesis/refs.bib.

\section*{The locked geometry panel, and why each piece is in it}

This thesis extracts five things from every coronary centerline: three curvature-based numbers and
two decisions about how those numbers, and the local curvature they come from, are computed and
kept. All five are locked in here because each has an evidenced or a mechanistically argued
connection to wall shear stress (WSS), the frictional drag blood exerts on the artery wall as it
flows past. That drag is not uniform along a vessel; it falls, or reverses over the heartbeat,
wherever the vessel bends or divides, and the endothelial cells lining the wall respond to
chronically low drag by turning inflamed and plaque-prone, which is where and why plaque starts
\cite{chatzizisis2007ess}. Tracing the flow through human coronary trees after death, Asakura and
Karino found almost every plaque in exactly two places: the inner wall of a bend and the outer wall
just past a bifurcation \cite{asakura1990flow}. The five measures below exist to locate those two
places without running a flow simulation for every patient.

\subsection*{1. Mean absolute curvature, the primary measure}

\textbf{Curvature} $\kappa$ describes how sharply a curve bends at a point: a tight bend has a small
radius of curvature $R$ and a large $\kappa = 1/R$, and a straight stretch has $\kappa = 0$. A
vessel's curvature therefore changes continuously along its length, so the first locked measure
takes its average over the whole vessel,
\begin{equation}
  \bar\kappa_a = \frac{1}{L}\int_0^L \lvert\kappa(s)\rvert\,\mathrm{d}s ,
  \label{eq:kappa-a}
\end{equation}
where $L$ is the vessel's length and $s$ the distance traveled along it. Every local bend adds
positively to this average, because the absolute value keeps a bend to the left from canceling a
bend to the right; the same quantity can be computed by summing the turning angle between every pair
of consecutive short steps along the vessel and dividing by $L$, a construction called the sum of
angles metric \cite{bullitt2003tortuosity}.

Kashyap and colleagues tested this measure, and three others, against a hemodynamic simulation of
127 people without coronary artery disease: for each person's left main bifurcation, they computed
the percentage of the vessel wall running below 0.4~Pa of time-averaged wall shear stress, a
simulation's stand-in for a chronically low-shear zone, and regressed each candidate tortuosity
measure against it \cite{kashyap2022tortuosity}. Mean absolute curvature $\bar\kappa_a$
(\Cref{eq:kappa-a}) explained the most variance of any measure they tried, $R^2 = 0.112$
($p < 0.001$). The classic \textbf{tortuosity index}, the ratio of the vessel's arc length to the
straight-line distance between its ends, explained essentially none, $R^2 \approx 0$
($p = 0.865$): a vessel can be long relative to that straight line either because it winds back and
forth or because it sweeps through one long, gentle curve, and the ratio cannot tell those two
apart. This is why $\bar\kappa_a$, not the tortuosity index, is this thesis's primary locked
measure, and why the tortuosity index is not carried forward as a feature at all.

\subsection*{2. Root-mean-square curvature, the sharp-bend companion}

A single sharp kink and a long, gentle arc can share the same total turning and so the same
$\bar\kappa_a$, yet a sharp kink disturbs flow more. The second locked measure weights sharper bends
more heavily by squaring curvature before averaging,
\begin{equation}
  \kappa_{\mathrm{rms}} = \left(\frac{1}{L}\int_0^L \kappa(s)^2\,\mathrm{d}s\right)^{1/2} .
  \label{eq:kappa-rms}
\end{equation}
In the same 127-person test, $\kappa_{\mathrm{rms}}$ also tracked the low-shear area,
$R^2 = 0.093$ ($p = 0.002$) \cite{kashyap2022tortuosity}, close to $\bar\kappa_a$ though not quite
as strongly. The same quantity, defined for retinal vessels as total squared curvature normalized
by length, was Hart and colleagues' original tortuosity measure, and separated tortuous from normal
retinal vessels with 91 to 95 percent accuracy \cite{hart1999tortuosity} (read at abstract level
only; the size of that classification test is not recorded here and should be checked before this
note is merged into the thesis). Because it squares a derivative, $\kappa_{\mathrm{rms}}$ is noisier
than $\bar\kappa_a$ and moves more with how much the centerline is smoothed before differentiating,
so it is reported only alongside a stated smoothing scale, never as a bare number.

\subsection*{3. Torsion and non-planarity, the out-of-plane part}

Curvature alone cannot tell a bend that stays flat from one that also twists out of its own plane,
the way a coiled spring twists as it bends. \textbf{Torsion} $\tau$ is the rate at which the plane
of the bend rotates about the vessel's own direction of travel: zero for a bend confined to one
plane, and nonzero once the vessel corkscrews out of it. Some simulation studies argue that this
out-of-plane twisting drives a spiral, helical flow pattern that protects the wall rather than
harming it \cite{shen2026geometry}, though that claim is read here at review level, from a table
summarizing other groups' simulations, not from the primary studies themselves. Because torsion
might help rather than hurt, it is kept as its own number rather than folded into $\bar\kappa_a$ or
$\kappa_{\mathrm{rms}}$, so a vessel that twists a great deal is never mistaken for one that simply
bends a great deal.

\subsection*{4. Bifurcation angle together with branch diameter}

At a bifurcation, the angle between the two daughter vessels changes where the incoming flow
separates from the wall, but the angle alone does not say how much flow each daughter carries: that
is set mostly by each daughter's own diameter. Pooling 1070 coronary trees, Taylor and colleagues
placed this flow-to-diameter relation, known as Murray's law, at an exponent of 2.39
(95\% CI 2.24 to 2.54), close to but below the cubic value classical theory predicts
\cite{taylor2024murray}. Because the drag a branch's wall feels depends steeply on how much flow it
carries, two bifurcations with the same angle but different diameter ratios load their walls
differently. The fourth locked measure is therefore the angle and the two daughter diameters
together, extracted as one unit, rather than the angle on its own.

\subsection*{5. Keeping curvature as a profile, not one number per vessel}

The first three measures each reduce a whole vessel to a single number, and a single number can hide
where along the vessel the bending actually sits. Tello Ayala and colleagues trained a model to read
the turning-angle profile along each vessel directly, rather than its average, on 22{,}334 patients
undergoing angiography, and it predicted coronary artery disease better than the same information
collapsed to its global mean, an area under the receiver operating characteristic curve (AUROC) of
0.67 against 0.60 for the mean alone \cite{telloayala2026tortuosity} (read at abstract level only).
The fifth locked choice follows from this: this thesis keeps the curvature profile $\kappa(s)$ along
each vessel, per segment, rather than collapsing it to $\bar\kappa_a$ immediately, so that where a
bend sits is never thrown away before it can be used.

\subsection*{Considered, not locked}

Three further ways of scoring tortuosity were considered and left out. Grisan and colleagues'
tortuosity density cuts a vessel at each point it changes bending direction and combines the
arc-to-chord ratio of every resulting piece \cite{grisan2008tortuosity}; it captures turning as well
as elongation, but it inherits whatever noise the inflection-point detection carries, on top of the
noise curvature already carries, for no clear gain over $\bar\kappa_a$ on the one evidenced test
above. Categorical clinical scores, such as counting bends sharper than 45 degrees or consecutive
180-degree turns \cite{groves2009tortuosity}, are simple to read off a scan but throw away the
continuous magnitude that $\bar\kappa_a$ and $\kappa_{\mathrm{rms}}$ keep, and neither has been
tested against a shear simulation the way the measures above have. And a single tortuosity score per
patient, rather than per vessel segment, was ruled out directly: across 1010 people undergoing
angiography, more tortuous arteries were, if anything, associated with less coronary artery disease,
not more \cite{li2011tortuosity}, the opposite of what a whole-patient score implies if disease
tracked tortuosity everywhere. Flow acts locally, so this thesis measures locally, one vessel
segment at a time.

None of the five locked measures is a direct measurement of wall shear stress, and a high score on
one is a candidate risk marker, not proof of risk. Even a full hemodynamic simulation only
moderately separated lesions that went on to cause a heart attack from those that did not, in
patients who already had one, area under the curve 0.63 to 0.68 \cite{griffo2026wss}; a shape
measure computed without running that simulation should not be expected to do better.
